% Options for packages loaded elsewhere
\PassOptionsToPackage{unicode}{hyperref}
\PassOptionsToPackage{hyphens}{url}
\documentclass[
  12pt,
]{article}
\usepackage{xcolor}
\usepackage[margin=1in]{geometry}
\usepackage{amsmath,amssymb}
\setcounter{secnumdepth}{-\maxdimen} % remove section numbering
\usepackage{iftex}
\ifPDFTeX
  \usepackage[T1]{fontenc}
  \usepackage[utf8]{inputenc}
  \usepackage{textcomp} % provide euro and other symbols
\else % if luatex or xetex
  \usepackage{unicode-math} % this also loads fontspec
  \defaultfontfeatures{Scale=MatchLowercase}
  \defaultfontfeatures[\rmfamily]{Ligatures=TeX,Scale=1}
\fi
\usepackage{lmodern}
\ifPDFTeX\else
  % xetex/luatex font selection
\fi
% Use upquote if available, for straight quotes in verbatim environments
\IfFileExists{upquote.sty}{\usepackage{upquote}}{}
\IfFileExists{microtype.sty}{% use microtype if available
  \usepackage[]{microtype}
  \UseMicrotypeSet[protrusion]{basicmath} % disable protrusion for tt fonts
}{}
\makeatletter
\@ifundefined{KOMAClassName}{% if non-KOMA class
  \IfFileExists{parskip.sty}{%
    \usepackage{parskip}
  }{% else
    \setlength{\parindent}{0pt}
    \setlength{\parskip}{6pt plus 2pt minus 1pt}}
}{% if KOMA class
  \KOMAoptions{parskip=half}}
\makeatother
\usepackage{longtable,booktabs,array}
\usepackage{calc} % for calculating minipage widths
% Correct order of tables after \paragraph or \subparagraph
\usepackage{etoolbox}
\makeatletter
\patchcmd\longtable{\par}{\if@noskipsec\mbox{}\fi\par}{}{}
\makeatother
% Allow footnotes in longtable head/foot
\IfFileExists{footnotehyper.sty}{\usepackage{footnotehyper}}{\usepackage{footnote}}
\makesavenoteenv{longtable}
% definitions for citeproc citations
\NewDocumentCommand\citeproctext{}{}
\NewDocumentCommand\citeproc{mm}{%
  \begingroup\def\citeproctext{#2}\cite{#1}\endgroup}
\makeatletter
 % allow citations to break across lines
 \let\@cite@ofmt\@firstofone
 % avoid brackets around text for \cite:
 \def\@biblabel#1{}
 \def\@cite#1#2{{#1\if@tempswa , #2\fi}}
\makeatother
\newlength{\cslhangindent}
\setlength{\cslhangindent}{1.5em}
\newlength{\csllabelwidth}
\setlength{\csllabelwidth}{3em}
\newenvironment{CSLReferences}[2] % #1 hanging-indent, #2 entry-spacing
 {\begin{list}{}{%
  \setlength{\itemindent}{0pt}
  \setlength{\leftmargin}{0pt}
  \setlength{\parsep}{0pt}
  % turn on hanging indent if param 1 is 1
  \ifodd #1
   \setlength{\leftmargin}{\cslhangindent}
   \setlength{\itemindent}{-1\cslhangindent}
  \fi
  % set entry spacing
  \setlength{\itemsep}{#2\baselineskip}}}
 {\end{list}}
\usepackage{calc}
\newcommand{\CSLBlock}[1]{\hfill\break\parbox[t]{\linewidth}{\strut\ignorespaces#1\strut}}
\newcommand{\CSLLeftMargin}[1]{\parbox[t]{\csllabelwidth}{\strut#1\strut}}
\newcommand{\CSLRightInline}[1]{\parbox[t]{\linewidth - \csllabelwidth}{\strut#1\strut}}
\newcommand{\CSLIndent}[1]{\hspace{\cslhangindent}#1}
\setlength{\emergencystretch}{3em} % prevent overfull lines
\providecommand{\tightlist}{%
  \setlength{\itemsep}{0pt}\setlength{\parskip}{0pt}}
\usepackage{bookmark}
\IfFileExists{xurl.sty}{\usepackage{xurl}}{} % add URL line breaks if available
\urlstyle{same}
\hypersetup{
  hidelinks,
  pdfcreator={LaTeX via pandoc}}

\author{}
\date{}

\begin{document}

\section{Buying Attention, Not Endorsement: Multimodal Clickbait and the
Bifurcation of User Engagement in Health-Science Videos on
Bilibili}\label{buying-attention-not-endorsement-multimodal-clickbait-and-the-bifurcation-of-user-engagement-in-health-science-videos-on-bilibili}

\begin{center}\rule{0.5\linewidth}{0.5pt}\end{center}

\subsection{Abstract}\label{abstract}

Short-video platforms have become a major entry point for public health
information, yet the visual cues that producers attach to thumbnails
remain undertheorized in health communication. This study examines
whether \emph{multimodal clickbait}---the use of dramatized facial
emotion, threat imagery, and dense text overlay on video thumbnails---is
associated with a distinct \emph{bifurcation} of user engagement on
Bilibili, China's largest long-form Chinese video platform. Drawing on
curiosity-gap, threat-arousal, and signaling perspectives (Scott 2021;
Lang 2000; Bird and Smith 2005; Donath 2007), we predict that multimodal
clickbait increases low-cost attention (views, likes) more than
high-cost endorsement (coins, favorites, shares), and that
thumbnail-title representativeness moderates this gap. We test these
predictions on a corpus of 4,562 unique health-science videos posted
between June 2016 and May 2026, collected via Bilibili's public search
API across 28 keywords spanning five health themes. Each video is paired
with its locally archived thumbnail, six interaction metrics, and
uploader-level attributes (follower count, platform verification,
original/reposted status). Visual constructs are coded using a
three-tier hybrid protocol: 500 thumbnails are double-blind coded by two
trained annotators against a five-dimension codebook, the full corpus is
processed by computer-vision (PaddleOCR, RetinaFace + FER) and a
vision-language model (Claude Opus 4.7), and convergent validity between
human and computational scores is established before scaling. We
estimate negative binomial and zero-inflated negative binomial models
with cluster-robust standard errors at the uploader level, and report
between-uploader and within-uploader fixed-effects specifications in
parallel. We pre-register five hypotheses (H1a--H4) and one exploratory
three-way interaction (H5). The design is observational; we do not make
causal claims and bound unobserved confounding using Oster δ and
E-values. This pre-results manuscript sets out the theoretical
motivation, measurement protocol, and analytic plan; empirical results
will be added after pre-registration on the Open Science Framework.
\emph{Keywords}: visual clickbait, health communication, Bilibili,
multimodal, user engagement, computational measurement.

\begin{center}\rule{0.5\linewidth}{0.5pt}\end{center}

\subsection{1. Introduction}\label{introduction}

A Bilibili user opening the platform's search page for ``癌症 早期信号''
(early signals of cancer) sees a vertical stack of video thumbnails
before reading a single full sentence of medical content. One thumbnail
shows a physician in a white coat behind a bold yellow caption reading
``千万别乱用，用错后果严重'' (do not misuse, the consequences are
severe) flanked by three red exclamation marks. Another shows a small
child crying next to an elderly man in a wheelchair, with a yellow
caption asking strangers not to ignore them. A third shows the same
physician set against a desktop with a clinical report, captioned only
``糖尿病出现口干口渴 有3种情况'' (when diabetes patients experience dry
mouth and thirst, there are three possibilities). All three appear in
the same search result. All three are tagged as health education by
their uploaders. Yet they ask the viewer to do very different things,
and they are likely to be rewarded in very different ways by the
platform's distinctive engagement architecture.

Bilibili distinguishes itself from short-video competitors by retaining
a layered, cost-graded set of user actions. Viewers can play
(\texttt{play}), like (\texttt{like}), post a danmaku floating comment
(\texttt{danmaku}), favorite for later retrieval (\texttt{favorites}),
share externally (\texttt{share}), and, crucially, spend a
platform-rationed token to ``coin'' the video (\texttt{coin}). Coins are
scarce: every user is allotted a small daily quota and they may give at
most one coin to any single video, so coining is an act that the
platform's design forces users to economize. In the broader literature
on platform engagement, this stratification is unusual. Most platforms
collapse engagement into views, likes, comments, and shares; few build
in an explicit, tokenized, cost-bearing form of endorsement (Lu and Shen
2023). Bilibili therefore offers a natural empirical environment in
which to ask whether different \emph{types} of engagement respond
differently to the same content cue.

The cue we focus on is the video thumbnail. Thumbnails are the primary
visual signal in a feed-based platform: they appear before any video
plays, occupy most of the user's field of view, and increasingly carry a
high density of designed elements---facial expressions, threat imagery,
color contrast, and overlaid text that essentially serves as a second,
parallel headline. A growing body of work in computational communication
treats thumbnails as standalone units of analysis (Lu and Pan 2022;
Al-Ali and Hamzeh 2024; Limpijankit and Kender 2025; Naveed, Uzmi, and
Qazi 2025). In journalism studies, ``clickbait'' has long denoted
headline strategies designed to maximize click-through at potential cost
to information quality or downstream trust (Munger 2020; Vultee et al.
2022; Shin, DeFelice, and Kim 2025; Wang et al. 2025). We extend this
conceptual move to the visual plane and ask: when health-science videos
on Bilibili use multimodal clickbait cues on their thumbnails, what kind
of user engagement do those cues actually produce?

Our argument has three parts. First, we treat Bilibili's engagement
metrics as forming a \emph{cost hierarchy}. Viewing and liking are
low-cost actions; coining, favoriting, and external sharing impose a
real budgetary or social cost on the user (Dong et al. 2025). Second,
drawing on curiosity-gap theory (Scott 2021; Blom and Hansen 2015), the
limited-capacity model of mediated message processing (Lang 2000),
visual framing research (Powell et al. 2015; Geise and Xu 2025),
experimental evidence that emotionally valenced images shape both
attention and selection on social-media news (Keib et al. 2018), and the
broader finding that clickbait headlines depress perceived credibility
even as they raise click rates (Molyneux and Coddington 2020; Khawar and
Boukes 2025), we expect multimodal clickbait to be more strongly
associated with low-cost attention than with high-cost endorsement.
Third, we expect this gap to be moderated by \emph{thumbnail--title
representativeness}: when the thumbnail accurately previews what the
title (and, by extension, the video) actually delivers, high-cost
endorsement is more likely to materialize, consistent with both the
image--text congruence literature (Li and Xie 2020; Cao, Li, and Zhang
2025; Yoon, Yoon, and Park 2024) and eye-tracking evidence on
safety-message attention (Klein et al. 2020).

We test these predictions on 4,562 unique health-science videos posted
on Bilibili between June 2016 and May 2026. The corpus was collected
through the platform's public search API across 28 Chinese-language
keywords spanning five health-content themes: general health education,
chronic disease, oncology and screening, lifestyle, and pre-defined
clickbait stem phrases (``medical doctor reminds you,'' ``doctors
finally admit,'' ``this habit causes cancer''). For each video we
collected complete API metadata, six interaction metrics, uploader-level
attributes including follower count and platform verification, and the
cover thumbnail itself. We code the visual constructs through a
three-tier hybrid protocol that combines double-blind human annotation
on a 500-thumbnail subsample, computer-vision pipelines for facial
emotion and text-overlay geometry, and vision-language model scoring on
the full corpus, with explicit convergent validity checks against the
human gold standard.

This study is observational and pre-results. We deliberately separate
the design and measurement portion of the project, presented here, from
the empirical estimation portion, which will follow Open Science
Framework pre-registration. Our purpose in the present paper is to (a)
lay out the theoretical case for treating visual clickbait as multimodal
and engagement as cost-graded; (b) document a measurement protocol that
audits hybrid human--machine coding in a way that subsequent work can
adopt; and (c) specify in advance an analytic plan that handles the
heavy zero inflation in the high-cost engagement variables, the
within-uploader clustering, and the well-known limits of associational
inference under recommender-system confounding (Lu and Pan 2021). We
organize the paper around four research questions:

\begin{itemize}
\tightlist
\item
  \textbf{RQ1}. To what extent do health-science videos on Bilibili use
  multimodal clickbait cues (dramatized emotion, threat imagery, dense
  text overlay) on their thumbnails?
\item
  \textbf{RQ2}. Is multimodal clickbait associated with higher low-cost
  engagement (views, likes) than with high-cost endorsement (coins,
  favorites, shares)?
\item
  \textbf{RQ3}. Does thumbnail--title representativeness moderate the
  gap between low-cost attention and high-cost endorsement?
\item
  \textbf{RQ4}. Does the bifurcation pattern depend on uploader
  verification status?
\end{itemize}

The remainder of the paper proceeds as follows. Section 2 reviews three
literatures we draw on: visual clickbait and thumbnail design,
image--text congruence and processing fluency, and platform-level
engagement stratification with particular attention to Bilibili. Section
3 develops our theoretical argument and formalizes five hypotheses.
Section 4 describes the corpus, the codebook, the hybrid annotation
pipeline, and the planned analysis, including a candid statement of the
causal boundary of the design.

\begin{center}\rule{0.5\linewidth}{0.5pt}\end{center}

\subsection{2. Literature Review}\label{literature-review}

\subsubsection{2.1 From clickbait headlines to multimodal clickbait
thumbnails}\label{from-clickbait-headlines-to-multimodal-clickbait-thumbnails}

Clickbait research originated in journalism studies as an attempt to
characterize headline strategies that withhold information in order to
maximize click-through. Linguistic analyses of online news identified
``forward-reference''---the use of cataphoric or deictic devices that
promise content the reader has not yet seen---as a signature device for
inducing anticipation (Blom and Hansen 2015). Relevance-theoretic work
extended this account by showing that successful clickbait headlines
exploit definite referring expressions and intensifiers to construct an
information gap that only the click can close (Scott 2021). The
economics of clickbait have been formalized as a strategic interaction
between an attention-constrained reader and a revenue-constrained
publisher, in which exaggeration is a rational response to a thin
attention market (Munger 2020). Large-scale corpus analyses of digital
newspaper headlines have isolated the specific linguistic
features---signal words, length, sentiment, numeric specificity---that
predict click-through in production environments (Kuiken et al. 2017),
and recent comparative work has documented that online-native outlets
use sensationalized headlines more heavily than legacy outlets, with
engagement consequences that differ across user segments (Khawar and
Boukes 2025). A parallel line of work in non-English settings shows that
the same headline mechanics travel to Chinese-language local media,
where length, question framing, and numerical claims predict clicks net
of content (Guo et al. 2026). Empirical work in this tradition has
further shown that clickbait headlines impose downstream costs on
perceptions of source credibility (Molyneux and Coddington 2020), on
subsequent attitudes toward the outlet (Vultee et al. 2022), and on
willingness to share content; emotional framing---particularly
rage-bait---operates differently from information-deficit framing (Shin,
DeFelice, and Kim 2025; Xu, Laffidy, and Ellis 2023). Reader-side
experimental work using EEG and behavioral measures finds that different
\emph{types} of clickbait (hyperbole, insinuation, visual rhetoric,
puzzle) elicit measurably different cognitive and affective responses
(Wang et al. 2025). The Persuasion Knowledge Model offers a broader lens
on why this matters: once readers recognize a persuasive intent, their
downstream evaluations of source and content change (Isaac 2025).

Two extensions of this body of work motivate the present study. The
first is the move to non-Western, platform-native, and visually dominant
settings. Lu and Pan, in their study of the Chinese government's
clickbait practices, demonstrated that clickbait is not merely a tabloid
phenomenon but a structural response to algorithmic visibility pressure
(Lu and Pan 2021). Their follow-up work on Douyin showed that visually
dominant short-video platforms have generated a distinctive grammar of
attention-maximizing video features, including high brightness, warm
color, short duration, and visual similarity to celebrity content (Lu
and Pan 2022). Their analysis of Chinese-language fact-checking videos
on Douyin further demonstrated that engagement and informational quality
respond to distinct multimodal features (Lu and Shen 2023). The second
extension is the move from text headlines to visual thumbnails. Al-Ali
and Hamzeh, analyzing Arabic-language clickbait thumbnails on YouTube,
showed that visual cues, overlaid text, and embedded punctuation jointly
constitute a meaning system that the text-only clickbait literature
systematically misses (Al-Ali and Hamzeh 2024). ThumbnailTruth, a recent
multi-modal LLM benchmark across cultures, formalized the detection
problem and reported that culturally diverse thumbnails resist
single-language solutions (Naveed, Uzmi, and Qazi 2025). Limpijankit and
Kender, using computational aesthetics, showed that systematic cultural
differences in thumbnail design can be measured at scale (Limpijankit
and Kender 2025). We build on these moves by treating \emph{visual
clickbait} as multimodal: a thumbnail is clickbait not only when its
overlaid text plays on a curiosity gap, but also when its facial
expressions are exaggerated, its threat imagery is salient, and its text
layer crowds out the underlying image.

\subsubsection{2.2 Image--text congruence, processing fluency, and
downstream
engagement}\label{imagetext-congruence-processing-fluency-and-downstream-engagement}

A second body of work, mostly in marketing and computational social
science, asks how the \emph{relation} between an image and its
accompanying text shapes user behavior. Li and Xie's analysis of branded
social-media posts established that image content, image quality, the
presence of human faces, and image--text matching each independently
predict engagement (Li and Xie 2020). Subsequent work using
deep-learning measures of image--text congruence has complicated the
simple ``more congruence is better'' intuition. Cao, Li, and Zhang
showed that the relationship between image--text congruence and consumer
preference can be non-monotonic: very high congruence delivers fluency,
very low congruence delivers surprise and elaboration, and the middle
ground may be the least attractive (Cao, Li, and Zhang 2025). In a
natural-language-processing reformulation, Yoon and colleagues argued
that news-thumbnail representativeness can be measured by counterfactual
text generation and that representativeness is a distinct construct from
headline--image similarity (Yoon, Yoon, and Park 2024). For visually
dominant platforms with high content volume, the question is no longer
whether congruence matters but \emph{which engagement metric} it matters
for. This is the empirical opening we exploit.

Visual framing research has documented that images and text contribute
non-redundantly to framing effects, with text shifting opinion and
images shifting behavioral intention (Powell et al. 2015). Eye-tracking
and selection experiments on social-media news further demonstrate that
emotionally valenced images alter not only what users \emph{attend to}
but also what they \emph{select to read} and \emph{share}, producing a
measurable wedge between initial attention and deeper engagement (Keib
et al. 2018). A recent systematic review of forty-five years of visual
framing studies found that the literature has paid disproportionate
attention to single-image stimuli and has lagged in multimodal,
platform-native, behavioral-outcome studies (Geise and Xu 2025). The
visual misinformation literature reaches a parallel conclusion: visual
misleading cues are systematically understudied (Heley, Gaysynsky, and
King 2022), and their effects on credibility and downstream behavior are
uneven across cultural and platform contexts (A. K. C. Liu and Kuru
2025). Eye-tracking evidence on health-safety messages specifically
shows that mismatch between an accompanying image and the textual
message degrades attention to and retention of the safety information
itself (Klein et al. 2020), suggesting that representativeness is not a
stylistic preference but a determinant of comprehension. In health
communication more broadly, scholarly calls have urged researchers to
move from descriptive visual analysis to outcome-oriented studies that
connect visual features to engagement and infodemic-relevant behavior
(King and Lazard 2020). Image-driven engagement on Twitter/X varies
systematically with the emotional content of climate imagery (Bravo et
al. 2025), and automated visual analysis is now a tractable empirical
strategy for studying social-media health effects (Peng, Lock, and Salah
2024; Araujo, Lock, and Velde 2020).

\subsubsection{2.3 Platform-level engagement stratification, Bilibili,
and
danmaku}\label{platform-level-engagement-stratification-bilibili-and-danmaku}

A third literature, focused on the Chinese-language video platform
Bilibili and its corporate cousin Douyin, supplies our outcome side.
Bilibili's distinctive feature is its dual register of viewer action:
alongside the standard view--like--comment--share suite, it preserves a
\emph{danmaku} layer of synchronously displayed comments and a
\emph{coin} token that the platform deliberately rations. The
fact-checking video literature has begun to take this stratification
seriously. Lu and Shen, analyzing Chinese-language fact-checking videos
on Douyin, distinguished engagement signals that respond to multimodal
features from those that do not (Lu and Shen 2023). On Bilibili itself,
Chen reported that danmaku interactions are predicted by content
features that differ from those predicting comment volume, suggesting
that danmaku occupies a behaviorally distinct slot in the engagement
hierarchy (Chen 2025). Dong and colleagues developed a typology of
danmaku ritual types (e.g., greeting, projection, evaluation) and showed
that ritual type, not raw danmaku volume, predicts downstream digital
engagement (Dong et al. 2025). These findings converge on a single
methodological point: collapsing engagement into a single score loses
theoretical information.

Signaling theory, originally developed in labor economics (Spence 1973)
and ported to online social environments by Donath (Donath 2007), offers
a complementary lens. Costly signals are informative precisely because
they are costly; cheap signals are common precisely because they are
cheap (Bird and Smith 2005). Applied to Bilibili, this implies that a
coin or a favorite carries information that a view or a like does not,
because the user has paid a real cost (a rationed token; the marginal
storage cost of a future-self bookmark) to issue it. The empirical
question is whether multimodal clickbait cues, which are designed to
maximize a cheap signal (clicks), do or do not also maximize the costly
signal. The Persuasion Knowledge Model predicts that once viewers
recognize persuasive intent, they will dampen costly endorsement even as
they continue to deliver cheap engagement (Isaac 2025). We do not test
that mechanism directly here, but it underlies our directional
predictions.

A separate strand of computational social science work documents how
scalable, multi-modal video analysis is now feasible at the scale of
tens of thousands of videos, but cautions that automated visual
measurement requires explicit construct-validity audits against human
annotation (Edelmann et al. 2020; Peng, Lock, and Salah 2024). This
methodological consensus shapes our annotation design (Section 4).

\subsubsection{2.4 What is missing}\label{what-is-missing}

The clickbait, image--text congruence, and Bilibili engagement
literatures have grown rapidly but have not yet been joined. Clickbait
research has remained largely text-centric; the few multimodal clickbait
studies are concentrated in English- and Arabic-language news settings
(Al-Ali and Hamzeh 2024; Naveed, Uzmi, and Qazi 2025). Image--text
congruence research is concentrated in branded marketing settings where
the outcome is purchase or attitude, not the engagement-cost ladder.
Bilibili engagement work is most developed at the danmaku layer and has
not, to our knowledge, addressed the question of whether different
visual cue types map differentially onto the platform's cost hierarchy.
Health-science content, with its distinctive stakes around credibility
and downstream behavior, has been studied for content quality but not
for the upstream visual cues that determine whether a viewer enters the
video at all. This study is positioned at the intersection of these
three gaps.

\begin{center}\rule{0.5\linewidth}{0.5pt}\end{center}

\subsection{3. Theory and Hypotheses}\label{theory-and-hypotheses}

\subsubsection{3.1 A cost-graded conception of
engagement}\label{a-cost-graded-conception-of-engagement}

The general framework for our predictions is the Differential
Susceptibility to Media Effects Model (DSMM), which holds that
media-effect estimates depend on the interaction of dispositional,
developmental, and social susceptibility with response states that
mediate the effect (Valkenburg and Peter 2013). We use DSMM as a
meta-frame: rather than predicting a single average effect of multimodal
clickbait, we predict a \emph{patterned} effect that differs across
engagement cost tiers because the dispositional cost of each action
shapes which response state (cheap attention vs.~costly endorsement)
dominates.

Our theoretical starting point is that user engagement on Bilibili is
not a single quantity but an ordered set of actions that differ in the
cost they impose on the user. \emph{View} costs the user only the
marginal attention required to begin a video. \emph{Like} costs the user
a tap. \emph{Danmaku} costs the user the time and self-disclosure
required to write and submit a public synchronous comment.
\emph{Favorite} costs the user the marginal cognitive overhead of
categorizing the video as worth re-watching. \emph{Share} costs the user
a social risk: the reputational implication of recommending the video to
an external audience. \emph{Coin}, finally, costs the user a fraction of
a daily platform-rationed budget and is irrevocable: a coin spent here
cannot be spent elsewhere today. We do not claim that all users perceive
these costs identically; we claim only that the ordering is robust
enough at the population level to motivate analytic separation.

This framing draws on signaling theory's core distinction between cheap
and costly signals (Spence 1973; Bird and Smith 2005; Donath 2007). The
information content of a signal in equilibrium is increasing in its
cost: cheap signals are common and uninformative; costly signals are
rare and diagnostic. Applied to a corpus of videos, this implies that
the same content feature can produce statistically distinct patterns in
cheap and costly engagement, and that pooling them obscures rather than
reveals the underlying behavior.

\subsubsection{3.2 Why multimodal clickbait should buy attention but not
endorsement}\label{why-multimodal-clickbait-should-buy-attention-but-not-endorsement}

Multimodal clickbait is a content strategy that maximizes the
probability of a click by amplifying perceptually salient cues in the
thumbnail. We identify three such cues: \emph{emotion intensity}
(exaggerated facial expressions, particularly fear, shock, or disgust);
\emph{threat imagery} (medical instruments, lesions, distressed bodies,
red warning marks); and \emph{text overlay intensity} (high-coverage,
high-contrast captions that pre-stage a curiosity gap). Each cue is
grounded in a distinct mechanism. The limited-capacity model of mediated
message processing holds that high-arousal, emotionally evocative
stimuli automatically allocate more cognitive resources to encoding,
increasing both initial attention and subsequent message recall (Lang
2000). Consistent with this account, emotion-laden imagery is processed
faster than neutral imagery, emotional intensity predicts virality of
social-media content (Bravo et al. 2025; Shin, DeFelice, and Kim 2025),
and experimental work has documented that emotionally valenced images in
news feeds shape attention, story selection, and downstream sharing
intent (Keib et al. 2018). Threat imagery activates the threat-salience
channel of fear appeals; classical health-behavior work formalizes this
channel within the Health Belief Model, in which perceived threat
operates jointly with perceived efficacy to produce action (Jones et al.
2015). On social media specifically, perceived threat and perceived
efficacy interact in a moderated-mediation path that elevates click
intention through fear arousal but dampens it when efficacy is low
(Zhang and Zhou 2019). Text-overlay intensity functions as a second
headline embedded in the image, multiplying surface area for
curiosity-gap construction (Al-Ali and Hamzeh 2024; Scott 2021).

A common feature of these mechanisms is that they target the
\emph{entry} decision---whether to click---rather than the
\emph{post-entry} decision of whether to endorse. Reader-side clickbait
studies are consistent with this asymmetry: clickbait headlines reliably
increase clicks but do not reliably increase trust, sharing, or other
downstream measures (Vultee et al. 2022; Shin, DeFelice, and Kim 2025;
Wang et al. 2025). The Persuasion Knowledge Model gives this asymmetry a
name: once viewers recognize that a thumbnail is engineered to maximize
their click probability, they may continue to click while withholding
the costly signals that would otherwise indicate endorsement (Isaac
2025). On Bilibili, the costly signals are precisely the rationed coin
and the future-oriented favorite and share. We therefore predict a
bifurcation:

\begin{itemize}
\tightlist
\item
  \textbf{H1a}. \emph{Multimodal clickbait intensity is positively
  associated with low-cost engagement (views, likes).}
\item
  \textbf{H1b}. \emph{Multimodal clickbait intensity is more weakly
  associated with high-cost endorsement (coins, favorites, shares) than
  with low-cost engagement.}
\end{itemize}

H1b is a \emph{comparative} hypothesis: it predicts that the
standardized coefficient on the multimodal clickbait index is smaller in
the high-cost equation than in the low-cost equation, not that it is
zero or negative. We test it with a non-parametric bootstrap of the
coefficient difference (Section 4.6).

\subsubsection{3.3 Why information gap should be treated
separately}\label{why-information-gap-should-be-treated-separately}

In our preliminary coding (Section 4.3) and consistent with the
curiosity-gap tradition, we distinguish \emph{information gap}---a
property of the title's promise rather than of the thumbnail's
image---from the three visual clickbait dimensions above. Information
gap is the extent to which the title and the visible thumbnail caption
deliberately withhold the punch line of the video. Examples in our
corpus include ``this habit causes cancer'' (which kind of habit?),
``doctors finally admit'' (admit what?), and ``if your physical exam
shows these signs, do not ignore them'' (which signs?). Linguistic work
on online news headlines identifies these constructions as deictic
forward-references: the headline uses a referring expression whose
antecedent is withheld, generating an anticipation that only the click
can resolve (Blom and Hansen 2015; Scott 2021). Theoretically, the
information-gap mechanism is cognitive rather than affective: it works
by making the viewer aware of a knowledge deficit and offering the click
as the closure. Because the gap is constructed in text, conflating it
with image-side cues such as facial emotion would misname the construct.
We therefore model information gap as a \emph{separate} predictor:

\begin{itemize}
\tightlist
\item
  \textbf{H1c}. \emph{Information-gap intensity is positively associated
  with low-cost engagement and more weakly associated with high-cost
  endorsement, in parallel to H1a--H1b.}
\end{itemize}

\subsubsection{3.4 Why thumbnail--title representativeness should
moderate the
gap}\label{why-thumbnailtitle-representativeness-should-moderate-the-gap}

Representativeness is the property of a thumbnail that it accurately
previews what the title---and by extension the video---is about. A
highly representative thumbnail of a video titled ``what 3 things to
know when diabetes patients experience dry mouth'' might display a
doctor with a clinical report and an overlay listing three items; an
unrepresentative thumbnail of the same video might display a generic
stock image of a glass of water. Two literatures lead to opposing
predictions, and the resolution is empirical.

The processing-fluency tradition holds that congruent image--text pairs
are easier to process, more credible-feeling, and more likely to convert
attention into action (Li and Xie 2020). Eye-tracking work in health
communication adds a behavioral mechanism: when an image and the
accompanying message do not match, viewers' attention shifts
inefficiently between the two modalities, and retention of the safety
information itself is degraded (Klein et al. 2020). The
curiosity-discrepancy tradition, by contrast, holds that mismatch
generates surprise and elaboration that can themselves be
conversion-enhancing under some conditions (Cao, Li, and Zhang 2025). We
expect the fluency channel to dominate for \emph{high-cost} endorsement,
because favoriting and coining require the viewer to make a
forward-looking judgment about content value, and that judgment depends
on whether the thumbnail's promise was kept. We expect
representativeness to have a weaker association with low-cost attention,
because attention is captured at the moment of the click, before any
representativeness comparison can be made:

\begin{itemize}
\tightlist
\item
  \textbf{H2}. \emph{Thumbnail--title representativeness is positively
  associated with high-cost endorsement.}
\item
  \textbf{H3}. \emph{The positive association between representativeness
  and high-cost endorsement is stronger than its association with
  low-cost engagement.}
\end{itemize}

\subsubsection{3.5 Medical authority cues and verification
status}\label{medical-authority-cues-and-verification-status}

Health-science content carries source-credibility implications that
general clickbait research does not. A white coat, a clinical report, or
a hospital backdrop function as visual signals of medical authority. The
same cues, however, are routinely co-opted by misleading health content;
recent commentary in \emph{Science Communication} has called for
systematic research on this dual-use property of visual health
information (Heley, Gaysynsky, and King 2022). Users themselves rely
heavily on social and heuristic shortcuts---rather than effortful
evaluation---when assessing online credibility, and visual authority
cues are exactly the kind of cognitively cheap, ``good-looking'' signal
that heuristic processing privileges (Metzger, Flanagin, and Medders
2010). In the Chinese-language digital news environment, recent
qualitative work documents that users articulate folk theories about
source trustworthiness that are sensitive to platform context and to
perceived sensationalism (S. Liu et al. 2023), reinforcing the
expectation that visual authority cues operate differently when paired
with high versus low clickbait intensity. The signaling-theory
expectation is that authority cues will amplify high-cost endorsement,
but only when they are \emph{consistent} with the underlying content;
when authority cues coexist with high clickbait intensity, viewers may
interpret the combination as packaging rather than credentialing,
dampening the endorsement response (Isaac 2025). We formalize the main
effect:

\begin{itemize}
\tightlist
\item
  \textbf{H4}. \emph{Medical authority visual cues are positively
  associated with high-cost endorsement.}
\end{itemize}

Bilibili distinguishes uploaders by formal verification. Verified
accounts (\texttt{is\_official\ ==\ TRUE}) include certified medical
professionals, institutional channels, and accredited media outlets. The
signaling and reader-side literatures both suggest that the bifurcation
in H1 should be \emph{attenuated} among verified accounts: viewers may
extend more trust to a verified uploader and convert clickbait-driven
attention into high-cost endorsement more readily; or, alternatively,
viewers may apply a stricter standard to a verified uploader and dampen
endorsement when verified accounts use clickbait. Because both
directions are theoretically defensible, we treat the moderation as
exploratory:

\begin{itemize}
\tightlist
\item
  \textbf{H5 (exploratory)}. \emph{The bifurcation pattern in H1 is
  moderated by uploader verification; we do not pre-commit to a
  direction.}
\end{itemize}

H5 is reported in the appendix only; the primary confirmatory family is
H1a--H4.

\subsubsection{3.6 What we are not
claiming}\label{what-we-are-not-claiming}

We are not claiming a causal effect of multimodal clickbait on
engagement. Our design is observational. Bilibili's recommender system
constitutes a substantial unobserved confounder of both the visual cues
that uploaders choose and the engagement that videos receive. We adopt
three protective measures: explicit associational language throughout,
sensitivity bounds (Oster δ and E-values) on the principal coefficients,
and a parallel within-uploader specification that absorbs all
time-invariant uploader characteristics, including the uploader's
tendency to be promoted by the recommender system. We do not claim that
even the within-uploader specification identifies a causal effect, only
that it strengthens the associational inference by partialing out a
large class of confounders.

\begin{center}\rule{0.5\linewidth}{0.5pt}\end{center}

\subsection{4. Data and Methods}\label{data-and-methods}

\subsubsection{4.1 Platform and corpus}\label{platform-and-corpus}

Bilibili (bilibili.com) is the largest long-form Chinese-language video
platform, founded in 2009 and best known for its synchronous danmaku
comment layer and its tokenized ``coin'' endorsement system. Unlike its
short-video competitors, Bilibili preserves long-form content alongside
short videos and retains an explicit four-way action set (\texttt{like},
\texttt{coin}, \texttt{favorite}, \texttt{share}) plus the
floating-comment \texttt{danmaku} channel. As of the third quarter of
2024, Bilibili reported an average of 348 million monthly active users
(Bilibili Inc. 2024). We choose Bilibili because its engagement
architecture, more than any other Chinese-language platform, separates
cheap and costly user actions at the design level.

We collected the corpus through the platform's public search API on May
24, 2026. The search API returns the same ranked list that the platform
serves to non-logged-in users entering a keyword query, with the
additional functionality of pagination and sort-order control. For each
of 28 Chinese-language keywords (Table 1), we issued three sort
queries---\texttt{totalrank} (relevance), \texttt{pubdate} (recency),
and \texttt{click} (popularity)---and crawled up to fifteen pages per
query. Keywords were organized into five thematic clusters: general
health education (e.g., 健康科普, 医生提醒, 体检报告; five keywords),
chronic disease (糖尿病, 高血压, 心脏病, 脂肪肝; four keywords),
oncology and screening (癌症 早期信号, HPV 疫苗, 肺癌 科普; three
keywords), lifestyle (减肥 科学, 脱发 医生, 睡眠 健康, 饮食健康; four
keywords), and pre-defined clickbait stem phrases (e.g., 医生提醒
千万别, 医生终于说了, 这个习惯 致癌; twelve keywords including
short-form variants). The clickbait-stem cluster was included
deliberately to ensure adequate variation in our principal independent
variable; we adjust for this by including keyword category fixed effects
in all primary models and by reporting a leave-one-keyword-out
robustness check.

{[}Table 1 about here{]}

For each returned video, we collected complete metadata (\texttt{bvid},
\texttt{aid}, title, description, publication timestamp, duration,
B-station category \texttt{tname}, original/reposted status
\texttt{copyright}), six interaction metrics (views, likes, coins,
favorites, shares, danmaku), uploader-level attributes (\texttt{mid},
follower count, platform verification status, channel signature,
uploader level), and the cover thumbnail (JPG, locally archived). After
deduplication on \texttt{bvid}, the working corpus contains 4,562 unique
videos from 2,731 unique uploaders (mean = 1.67 videos per uploader;
median = 1; maximum = 47). The publication dates span June 22, 2016 to
May 24, 2026, with 71 percent of videos published in 2024--2026. All
interaction metrics represent a single cross-sectional snapshot taken
within a six-hour window on May 24, 2026; we therefore introduce
\texttt{video\_age\_days} (the number of days between publication and
crawl) as a primary control. Thumbnail download succeeded for 4,564
images, providing 100 percent coverage of analyzable videos.

\subsubsection{4.2 Outcome variables and their
distributions}\label{outcome-variables-and-their-distributions}

We treat the six interaction metrics as an ordered set of engagement
actions stratified by cost. The full descriptive distribution is
reported in Table 2. Three features of the distribution are decisive for
the modeling strategy.

{[}Table 2 about here{]}

First, all six outcomes are heavily right-skewed, with maxima orders of
magnitude above their medians. Second, the rate of structural zeros
differs sharply across outcomes: 1.3 percent for views, 5.1 percent for
likes, 16.6 percent for favorites, 27.0 percent for shares, 30.5 percent
for comments, 36.5 percent for coins, and 45.8 percent for danmaku.
Third, the high-cost outcomes (coins, shares, danmaku) all exhibit
zero-inflation rates above 25 percent, indicating that a non-trivial
share of videos generate genuinely no costly endorsement. We accordingly
fit negative binomial (NB) models to the low-zero outcomes (views,
likes, favorites) and zero-inflated negative binomial (ZINB) models to
the high-zero outcomes (coins, shares, danmaku). The inflation equation
in the ZINB specification includes \texttt{log\_followers} and
\texttt{video\_age\_days}, which are theoretically the variables most
likely to predict structural absence of any endorsement.

\subsubsection{4.3 Codebook for visual
constructs}\label{codebook-for-visual-constructs}

Five visual constructs are coded per thumbnail: a binary topic-relevance
gate, three image-side clickbait dimensions, one text-side curiosity-gap
dimension, and a four-component medical-authority composite. The
complete codebook with operational definitions, exemplar images,
calibration rules, and adjudication procedures is included as
Supplementary Material
(\texttt{Data\ process/annotation/codebook\_visual\_clickbait.md}). We
summarize the constructs here.

\textbf{Topic relevance} (binary, gate variable). Coded \texttt{1} if
the video is recognizably a health-education attempt directed at a
general audience, \texttt{0} otherwise. The gate excludes emotional
appeals (e.g., medical-crowdfunding videos that use health imagery but
are not educational), fitness routines mislabeled as health science, and
entertainment content that uses health keywords for retrieval but does
not deliver health information. Videos with
\texttt{topic\_relevance\ ==\ 0} are excluded from all primary models.

\textbf{Emotion intensity} (ordinal, 0--3). Codes the degree of
exaggerated affective display on visible human faces in the thumbnail: 0
= neutral or no face; 1 = mild expression; 2 = pronounced surprise,
concern, or fear; 3 = extreme dramatized affect (open-mouth shock,
exaggerated grimaces, staged crying). The construct is grounded in the
high-arousal literature on viral content (Bravo et al. 2025; Shin,
DeFelice, and Kim 2025).

\textbf{Threat imagery} (ordinal, 0--2). Codes the salience of visual
threat cues: 0 = none; 1 = mild medical-setting cues (white coats,
clinical instruments, hospital backdrops) without explicit threat; 2 =
explicit threat cues (lesions, distressed bodies, red warning marks,
dramatic before/after images). The construct is grounded in the
visual-misinformation literature (Heley, Gaysynsky, and King 2022; A. K.
C. Liu and Kuru 2025) and in social-media fear-appeal research that has
documented how perceived threat shapes click intention (Zhang and Zhou
2019).

\textbf{Text overlay intensity} (ordinal, 0--3). Codes the proportion of
thumbnail area, font contrast, and emphasis density of overlaid text: 0
= no overlay; 1 = minimal caption; 2 = mid-density (a single short
caption with limited emphasis); 3 = high-density (large captions
covering 30 percent or more of the thumbnail area, often with
color-block backgrounds and exclamation marks). The construct is
grounded in the multimodal clickbait literature (Al-Ali and Hamzeh 2024;
Naveed, Uzmi, and Qazi 2025) and in evidence that image--text mismatch
alters cross-modal attention allocation (Klein et al. 2020).

\textbf{Information gap} (ordinal, 0--2). Codes the extent to which the
title and any thumbnail caption deliberately withhold the substantive
punch line of the video: 0 = no withholding (the title fully describes
the content); 1 = partial withholding (the title identifies the topic
but withholds the resolution); 2 = explicit curiosity gap (the title
constructs a ``this thing,'' ``these signs,'' ``what really happened''
puzzle). Operationalization follows the forward-reference and
relevance-theoretic accounts of clickbait headlines (Blom and Hansen
2015; Scott 2021). This construct is text-side, not image-side, and is
modeled separately from the visual clickbait dimensions.

\textbf{Medical authority cues} (4 binary items summed, 0--4). The four
items are: presence of a white coat or scrubs; presence of clinical
instruments (stethoscope, otoscope, syringe); presence of a clinical
report, scan, or chart; presence of an institutional logo or identifying
credential. The construct is grounded in source-credibility theory and
in the visual-misinformation literature on authority signaling (Isaac
2025; A. K. C. Liu and Kuru 2025).

\subsubsection{4.4 Hybrid annotation
pipeline}\label{hybrid-annotation-pipeline}

Visual constructs are scored using a three-tier hybrid protocol that
combines human gold-standard annotation, computer-vision feature
extraction, and vision-language model coding. The protocol is designed
to balance the construct-validity advantages of human coding against the
scale advantages of automated measurement, and to make the trade-off
legible to readers.

\textbf{Tier 1: Human gold standard (N = 500).} We draw a stratified
random subsample of 500 thumbnails from the 4,562-video corpus,
stratifying by keyword cluster (5 strata), uploader verification (2
strata), and the within-stratum tercile of view count (3 strata), to
ensure that all theoretically important regions of the corpus are
represented in the gold-standard set. Two trained annotators code each
thumbnail on all five constructs independently and blind to each other's
coding. Inter-rater reliability is computed as quadratic-weighted
Cohen's κ for the four ordinal constructs and as unweighted Cohen's κ
for the binary topic-relevance gate, with Krippendorff's α reported as a
robustness statistic. Target IRR is κ ≥ 0.85 for the gate variable and κ
≥ 0.70 for the four substantive constructs. Pairs with disagreement of ≥
2 ordinal steps are adjudicated by the principal investigator.
Annotators receive a fifteen-page codebook, a twenty-thumbnail
calibration set, and a one-hour group calibration session before
independent coding begins.

\textbf{Tier 2: Full-corpus computational coding.} All 4,562 thumbnails
are processed through two parallel pipelines. The computer-vision
pipeline follows the Automated Visual Content Analysis protocol
developed for communication research (Araujo, Lock, and Velde 2020): we
use RapidOCR for text detection and geometric analysis (yielding
text-area coverage and font-contrast measures that feed the
\texttt{text\_overlay\_intensity} score) and RetinaFace + a
Chinese-validated FER model for facial-emotion recognition (yielding a
continuous emotion-intensity score). The vision-language pipeline issues
a single structured prompt to Claude Opus 4.7 (Anthropic;
\texttt{claude-opus-4-7}) with the codebook operational definitions and
a small set of in-prompt exemplars, returning ordinal scores on
\texttt{topic\_relevance}, \texttt{threat\_imagery}, \texttt{info\_gap},
and the four medical-authority items. The raw VLM API responses are
archived for replication. For thumbnail--title representativeness, we
use Chinese-CLIP (ViT-B/16) (Yang et al. 2022) to compute cosine
similarity between cover-image and title embeddings, validated against a
small ordinal human rating on the 500-thumbnail subsample.

\textbf{Tier 3: Convergent-validity audit.} Computational scores are
validated against the human gold standard on the 500-thumbnail subsample
using Spearman's ρ. We pre-specify a threshold of ρ ≥ 0.60 for a
computational measure to be used in the full-sample analysis. Measures
that fall short of this threshold are reported only on the 500-thumbnail
gold-standard subsample. Computational measures that exceed the
threshold are used at full corpus scale, with all primary models
additionally fit on the 500-thumbnail gold-standard subsample as a
robustness check.

This pipeline is broadly consistent with the methodological consensus in
computational visual social science (Edelmann et al. 2020; Peng, Lock,
and Salah 2024) and with the multimodal annotation strategies adopted by
recent thumbnail-detection benchmarks (Naveed, Uzmi, and Qazi 2025). The
departure from prior practice is the explicit, pre-registered Spearman
threshold for promoting a computational measure to full-corpus use.

\subsubsection{4.5 Control variables}\label{control-variables}

The principal control variables are: \texttt{video\_age\_days}, the
number of days between publication and crawl, included to absorb the
cumulative-engagement bias toward older videos; \texttt{log\_followers},
the log-transformed follower count of the uploader, included to absorb
channel size; \texttt{duration\_log}, the log-transformed duration in
seconds; \texttt{pubdate\_year}, modeled as a categorical fixed effect
to absorb platform-mechanism evolution across the ten-year publication
window; \texttt{tname}, the B-station category fixed effect;
\texttt{copyright\_self}, a binary indicator that the video is original
(as opposed to reposted, which is rare in our sample but theoretically
distinct); and the binary \texttt{is\_official} verification indicator.
We deliberately exclude \texttt{order\_source} (the sort order under
which the video was retrieved) from the control set, because it is
endogenous to the video's clickbait intensity and thus post-treatment.
We report a sensitivity model that includes \texttt{order\_source} fixed
effects in the appendix. For the 13 percent of videos with missing
values on \texttt{tid}, \texttt{copyright}, or \texttt{uploader\_level}
(a legacy of an earlier crawl batch), we use isotonic-regression
imputation of \texttt{video\_age\_days} from the monotonic \texttt{aid}
index and code the remaining missing fields as a separate ``missing''
category.

\subsubsection{4.6 Analytic strategy}\label{analytic-strategy}

For each of the six outcomes, we fit a sequence of five nested
specifications. Model M1 includes the multimodal clickbait index and the
information-gap index alone. Model M2 adds the full control set. Model
M3 adds the representativeness main effect. Model M4 adds the
multimodal-clickbait-by-representativeness and
information-gap-by-representativeness interactions. Model M5, the
exploratory specification, additionally tests the three-way interaction
with \texttt{is\_official}. Standard errors are cluster-robust at the
uploader (\texttt{mid}) level to account for non-independence among
videos posted by the same channel.

We pre-register a \emph{dual-track} primary specification: a
between-uploader cross-sectional model (M2a) on the full 4,562-video
sample, and a within-uploader fixed-effects model (M2b) on the subsample
of uploaders with ≥ 2 videos (n ≈ 1,400 uploaders contributing ≈ 3,500
videos). The two-track design is motivated by the simultaneous threat of
unobserved confounding from the recommender system at the
between-uploader level and the loss of statistical power at the
within-uploader level. We report both as primary and treat their
agreement (or disagreement) as substantive.

To test the comparative hypothesis H1b, we estimate the bifurcation
contrast as the difference between the standardized multimodal-clickbait
coefficient in the low-cost equation (views; NB) and in the high-cost
equation (coins; ZINB). We construct a 95 percent confidence interval
for this difference using a non-parametric bootstrap on 1,000 resamples
of the video-level data, stratified by uploader. All marginal effects
are reported as average marginal effects (AMEs) using the
\texttt{marginaleffects} R package (Arel-Bundock, Greifer, and Heiss
2024), consistent with the journal's reporting norms.

We pre-specify multiple-comparisons control as follows: the family of
confirmatory tests is H1a, H1b, H1c, H2, H3, and H4; for this family, we
apply Benjamini--Hochberg false-discovery-rate control at q = 0.05, with
the high-cost coin and favorites outcomes designated as primary and the
remaining four as secondary. The exploratory H5 is reported without
multiple-comparisons adjustment.

\subsubsection{4.7 Robustness battery}\label{robustness-battery}

We pre-register the following sensitivity checks: (R1) re-fitting with
ratio outcomes (coins/views, favorites/views, shares/views) using beta
regression; (R2) restricting to the post-2022 subsample, to address
platform-mechanism evolution; (R3) trimming the top 5 percent of
channels by follower count; (R4) restricting to the ``Science \&
Knowledge'' tname category; (R5) Poisson rather than NB; (R6) replacing
the continuous multimodal-clickbait index with quartile dummies, to
relax the linearity assumption; (R7) Oster's δ for
unobserved-confounding bound (Oster 2019); (R8) E-values for
treatment-effect robustness (VanderWeele and Ding 2017); (R9)
restriction to videos with \texttt{topic\_relevance\ ==\ 1} under the
strict human gold-standard coding; (R10) human-coded-only on the
500-thumbnail subsample vs.~VLM-coded on the full corpus, to detect
measurement-source dependence; (R11) listwise deletion vs.~multiple
imputation (m = 20) for the 13 percent missingness; (R12) two-way
clustering on (\texttt{mid}, \texttt{pubdate\_month}) to address
recommender-system-induced SUTVA violations; and (R13)
leave-one-keyword-out across the five clickbait-stem seed keywords.

\subsubsection{4.8 Causal boundary and
limitations}\label{causal-boundary-and-limitations}

We do not claim a causal effect of multimodal clickbait on user
engagement. Bilibili's recommender system jointly determines both which
thumbnails get promoted into the feed and which receive subsequent
engagement; the backdoor between the thumbnail and the engagement metric
cannot be closed with observational data alone. The within-uploader
fixed-effects specification absorbs all time-invariant uploader
characteristics, including a channel's average historical promotion by
the recommender, but does not absorb video-level promotion shocks. Oster
δ and E-value sensitivity analyses quantify the magnitude of unobserved
confounding that would be required to nullify the estimated
associations, but they are not a substitute for an exogenous source of
variation. Three further limitations apply. First, our corpus is
``videos that were returned by the search API for 28 health-related
keywords,'' not ``all health-related videos on Bilibili''; results
generalize only to the search-indexed subset. Second, engagement metrics
are a single-day snapshot; ratio outcomes (R1) partially mitigate
cumulative-engagement bias but do not eliminate it. Third, we measure
visual features of the cover thumbnail but not the audio-visual content
of the video itself; a thumbnail-misleading video may be uniformly
clickbait-y or may deliver substantive content once entered, and our
design does not distinguish these.

\subsubsection{4.9 Pre-registration, open science, and
ethics}\label{pre-registration-open-science-and-ethics}

The pre-analysis plan, the codebook, the full set of VLM prompts, and
the analysis code will be deposited on the Open Science Framework prior
to the analysis of any outcome model. The corpus metadata (excluding raw
thumbnail images, which are subject to platform terms of service) will
be archived on Zenodo with a digital object identifier. Thumbnail images
themselves will not be redistributed; image URLs and a re-collection
script will be provided. The uploader identifier (\texttt{mid}) is
one-way hashed in the public dataset; the uploader signature field
(\texttt{uploader\_sign}), which may contain identifying titles, will be
redacted using a structured-label scheme. The study uses only public,
non-account-restricted platform metadata, with no interaction with users
or uploaders; our institution's review board has classified the protocol
as exempt under Category 4 (research involving the collection or study
of existing publicly available data). We disclose the use of Claude Opus
4.7 as a secondary coder in the annotation pipeline; all primary coding
decisions are retained by the human annotators, and the model is used
only for full-corpus scaling against the human gold standard.

\begin{center}\rule{0.5\linewidth}{0.5pt}\end{center}

\subsection{References}\label{references}

Al-Ali, M.N. and S.M. Hamzeh. 2024. ``Extra Cues Extra Views: A
Multimodal Detection of Arabic Clickbait Thumbnail Verbo-Visual Cues.''
\emph{Discourse \& Communication}.

Arel-Bundock, V., N. Greifer, and A. Heiss. 2024. ``How to Interpret
Statistical Models Using marginaleffects for R and Python.''
\emph{Journal of Statistical Software} 111(9): 1--32.
https://doi.org/10.18637/jss.v111.i09.

Araujo, T., I. Lock, and B. van de Velde. 2020. ``Automated Visual
Content Analysis (AVCA) in Communication Research: A Protocol for Large
Scale Image Classification with Pre-Trained Computer Vision Models.''
\emph{Communication Methods and Measures} 14(4): 239--265.

Bilibili Inc.~2024. \emph{Bilibili Inc.~Reports Third Quarter 2024
Financial Results}. Form 6-K filed with the U.S. Securities and Exchange
Commission, November 14, 2024. Retrieved May 25, 2026
(https://www.sec.gov/cgi-bin/browse-edgar?action=getcompany\&CIK=0001723690).

Bird, R.B. and E.A. Smith. 2005. ``Signaling Theory, Strategic
Interaction, and Symbolic Capital.'' \emph{Current Anthropology}.

Blom, J.N. and K.R. Hansen. 2015. ``Click Bait: Forward-Reference as
Lure in Online News Headlines.'' \emph{Journal of Pragmatics} 76:
87--100.

Bravo, I., K. Prasse, S. Walter, S. O'Neill, and M. Keuper. 2025.
``Global Dynamics of Climate Change Imagery: Emotional and Engagement
Effects Across Visual Frames on Twitter/X.'' \emph{Science
Communication}.

Cao, J., X. Li, and L. Zhang. 2025. ``Is Relevancy Everything? A
Deep-Learning Approach to Understand the Effect of Image-Text
Congruence.'' \emph{Management Science}.

Chen, J. 2025. ``Understanding Danmaku and Comment Interactions Through
Content Features and Video Popularity.'' \emph{Procedia Computer
Science}.

Donath, J. 2007. ``Signals in Social Supernets.'' \emph{Journal of
Computer-Mediated Communication}.

Dong, W., Y. Liu, W. Wang, L. Jiang, and Y. Yi. 2025. ``The Impact of
Danmaku Ritual Types on User Digital Engagement in Video-Based Social
Media: The Moderating Role of Influencer Types and Domains.''
\emph{Psychology \& Marketing}.

Edelmann, A., T. Wolff, D. Montagne, and C.A. Bail. 2020.
``Computational Social Science and Sociology.'' \emph{Annual Review of
Sociology}.

Geise, S. and Y. Xu. 2025. ``Effects of Visual Framing in Multimodal
Media Environments: A Systematic Review of Studies Between 1979 and
2023.'' \emph{Journalism \& Mass Communication Quarterly}.

Guo, L., Y. Wang, P. Li, Y. Wang, and Y. Li. 2026. ``The Impact of
Headline Characteristics on Clicks: A Case Study of a Chinese Local
Medium.'' \emph{Journalism Practice} 20(4): 1427--1455.

Heley, K., A. Gaysynsky, and A.J. King. 2022. ``Missing the Bigger
Picture: The Need for More Research on Visual Health Misinformation.''
\emph{Science Communication} 44(4): 514--527.

Isaac, M.S. 2025. ``Thirty Years of Persuasion Knowledge Research: From
Demonstrating Effects to Building Theory to Increasing Applicability.''
\emph{Consumer Psychology Review}.

Jones, C.L., J.D. Jensen, C.L. Scherr, N.R. Brown, K. Christy, and J.
Weaver. 2015. ``The Health Belief Model as an Explanatory Framework in
Communication Research: Exploring Parallel, Serial, and Moderated
Mediation.'' \emph{Health Communication} 30(6): 566--576.

Keib, K., C. Espina, Y.-I. Lee, B.W. Wojdynski, D. Choi, and H. Bang.
2018. ``Picture This: The Influence of Emotionally Valenced Images on
Attention, Selection, and Sharing of Social Media News.'' \emph{Media
Psychology} 21(2): 202--221.

Khawar, S. and M. Boukes. 2025. ``Analyzing Sensationalism in News on
Twitter (X): Clickbait Journalism by Legacy vs.~Online-Native Outlets
and the Consequences for User Engagement.'' \emph{Digital Journalism}
13(8): 1482--1502.

King, A.J. and A.J. Lazard. 2020. ``Advancing Visual Health
Communication Research to Improve Infodemic Response.'' \emph{Health
Communication} 35(14): 1723--1728.

Klein, E.G., K. Roberts, J. Manganello, R. McAdams, and L. McKenzie.
2020. ``When Social Media Images and Messages Don't Match: Attention to
Text versus Imagery to Effectively Convey Safety Information on Social
Media.'' \emph{Journal of Health Communication} 25(11): 879--884.

Kuiken, J., A. Schuth, M. Spitters, and M. Marx. 2017. ``Effective
Headlines of Newspaper Articles in a Digital Environment.''
\emph{Digital Journalism} 5(10): 1300--1314.

Lang, A. 2000. ``The Limited Capacity Model of Mediated Message
Processing.'' \emph{Journal of Communication} 50(1): 46--70.

Li, Y. and Y. Xie. 2020. ``Is a Picture Worth a Thousand Words? An
Empirical Study of Image Content and Social Media Engagement.''
\emph{Journal of Marketing Research}.

Limpijankit, M. and J.R. Kender. 2025. ``Detecting Cultural Differences
in News Video Thumbnails via Computational Aesthetics.''

Liu, A.K.C. and O. Kuru. 2025. ``Understanding the Effects of Visual
Misinformation: A Systematic Review of 10 Years (2014--2024).''
\emph{Mass Communication and Society}.

Liu, S., J. Xu, Z. Zhao, and X. Li. 2023. ``Factors Affecting Trust in
Chinese Digital Journalism: Approach Based on Folk Theories.''
\emph{Media and Communication} 11(4): 355--366.

Lu, Y. and J. Pan. 2021. ``Capturing Clicks: How the Chinese Government
Uses Clickbait to Compete for Visibility.'' \emph{Political
Communication}.

Lu, Y. and J. Pan. 2022. ``The Pervasive Presence of Chinese Government
Content on Douyin Trending Videos.'' \emph{Computational Communication
Research}.

Lu, Y. and C. Shen. 2023. ``Unpacking Multimodal Fact-Checking: Features
and Engagement of Fact-Checking Videos on Chinese TikTok (Douyin).''
\emph{Social Media + Society}.

Metzger, M.J., A.J. Flanagin, and R.B. Medders. 2010. ``Social and
Heuristic Approaches to Credibility Evaluation Online.'' \emph{Journal
of Communication} 60(3): 413--439.

Molyneux, L. and M. Coddington. 2020. ``Aggregation, Clickbait and Their
Effect on Perceptions of Journalistic Credibility and Quality.''
\emph{Journalism Practice} 14(4): 429--446.

Munger, K. 2020. ``All the News That's Fit to Click: The Economics of
Clickbait Media.'' \emph{Political Communication}.

Naveed, W., Z.A. Uzmi, and Z.A. Qazi. 2025. ``ThumbnailTruth: A
Multi-Modal LLM Approach for Detecting Misleading YouTube Thumbnails
Across Diverse Cultural Settings.''

Oster, E. 2019. ``Unobservable Selection and Coefficient Stability:
Theory and Evidence.'' \emph{Journal of Business \& Economic Statistics}
37(2): 187--204. https://doi.org/10.1080/07350015.2016.1227711.

Peng, Y., I. Lock, and A.A. Salah. 2024. ``Automated Visual Analysis for
the Study of Social Media Effects: Opportunities, Approaches, and
Challenges.'' \emph{Communication Methods and Measures}.

Powell, T.E., H.G. Boomgaarden, K. De Swert, and C.H. de Vreese. 2015.
``A Clearer Picture: The Contribution of Visuals and Text to Framing
Effects.'' \emph{Journal of Communication}.

Scott, K. 2021. ``You Won't Believe What's in This Paper! Clickbait,
Relevance and the Curiosity Gap.'' \emph{Journal of Pragmatics} 175:
53--66.

Shin, J., C. DeFelice, and S. Kim. 2025. ``Emotion Sells: Rage Bait
vs.~Information Bait in Clickbait News Headlines on Social Media.''
\emph{Digital Journalism}.

Spence, M. 1973. ``Job Market Signaling.'' \emph{Quarterly Journal of
Economics}.

Valkenburg, P.M. and J. Peter. 2013. ``The Differential Susceptibility
to Media Effects Model.'' \emph{Journal of Communication} 63(2):
221--243.

VanderWeele, T.J. and P. Ding. 2017. ``Sensitivity Analysis in
Observational Research: Introducing the E-Value.'' \emph{Annals of
Internal Medicine} 167(4): 268--274. https://doi.org/10.7326/M16-2607.

Vultee, F., G.S. Burgess, D. Frazier, and K. Mesmer. 2022. ``Here's What
to Know About Clickbait: Effects of Image, Headline and Editing on
Audience Attitudes.'' \emph{Journalism Practice}.

Wang, Y., B. Hu, C. Tang, and X. Yang. 2025. ``Decoding Clickbait: The
Impact of Clickbait Types and Structures on Cognitive and Emotional
Responses in Online Interactions.'' \emph{Cyberpsychology, Behavior \&
Social Networking}.

Xu, Z., M. Laffidy, and L. Ellis. 2023. ``Clickbait for Climate Change:
Comparing Emotions in Headlines and Full-Texts and Their Engagement.''
\emph{Information, Communication \& Society} 26(10): 1915--1932.

Yang, A., J. Pan, J. Lin, R. Men, Y. Zhang, J. Zhou, and C. Zhou. 2022.
``Chinese CLIP: Contrastive Vision-Language Pretraining in Chinese.''

Yoon, Y., S. Yoon, and K. Park. 2024. ``Assessing News Thumbnail
Representativeness: Counterfactual Text Can Enhance the Cross-Modal
Matching Ability.'' \emph{Findings of the Association for Computational
Linguistics: ACL 2024}.

Zhang, X. and S. Zhou. 2019. ``Clicking Health Risk Messages on Social
Media: Moderated Mediation Paths Through Perceived Threat, Perceived
Efficacy, and Fear Arousal.'' \emph{Health Communication} 34(11):
1359--1368.

\begin{center}\rule{0.5\linewidth}{0.5pt}\end{center}

\subsection{Table 1: Keyword Strata and Per-Keyword Sample
Sizes}\label{table-1-keyword-strata-and-per-keyword-sample-sizes}

\begin{longtable}[]{@{}lll@{}}
\toprule\noalign{}
Cluster & Keyword & N \\
\midrule\noalign{}
\endhead
\bottomrule\noalign{}
\endlastfoot
General health education & 医学科普 & 435 \\
General health education & 健康科普 & 429 \\
General health education & 医生提醒 & 390 \\
General health education & 体检报告 & 248 \\
General health education & 饮食健康 & 196 \\
General health education & 养生科普 & 180 \\
Chronic disease & 糖尿病 科普 & 180 \\
Chronic disease & 高血压 科普 & 180 \\
Chronic disease & 心脏病 科普 & 180 \\
Chronic disease & 脂肪肝 科普 & 180 \\
Chronic disease & 糖尿病 (short) & 60 \\
Chronic disease & 高血压 (short) & 20 \\
Chronic disease & 高血压 危险信号 & 6 \\
Oncology and screening & 癌症 早期信号 & 180 \\
Oncology and screening & HPV 疫苗 & 180 \\
Oncology and screening & 肺癌 科普 & 180 \\
Oncology and screening & HPV (short) & 60 \\
Lifestyle & 减肥 科学 & 180 \\
Lifestyle & 脱发 医生 & 180 \\
Lifestyle & 睡眠 健康 & 180 \\
Lifestyle & 减肥 (short) & 60 \\
Lifestyle & 脱发 (short) & 60 \\
Clickbait stem & 医生提醒 千万别 & 180 \\
Clickbait stem & 医生终于说了 & 180 \\
Clickbait stem & 体检报告 异常 & 130 \\
Clickbait stem & 这个习惯 致癌 & 92 \\
Clickbait stem & 健康科普 千万别吃 & 18 \\
Clickbait stem & 医生提醒 致癌 & 18 \\
\textbf{Total (pre-dedup)} & --- & \textbf{5,062} \\
\textbf{Total (post-dedup, primary corpus)} & --- & \textbf{4,562} \\
\end{longtable}

\textbf{Notes}: N reflects the number of unique videos retrieved under
each keyword × sort-order combination, before global deduplication on
\texttt{bvid}. Post-deduplication, 4,562 unique videos remain in the
analytic corpus. Keyword clusters are coded as a categorical fixed
effect in all primary models. The ``clickbait stem'' cluster includes
pre-defined search terms intended to oversample videos with high
multimodal clickbait intensity; a leave-one-keyword-out robustness check
is reported in Section 4.7.

\begin{center}\rule{0.5\linewidth}{0.5pt}\end{center}

\subsection{Table 2: Descriptive Distributions of Engagement Outcomes (N
=
4,562)}\label{table-2-descriptive-distributions-of-engagement-outcomes-n-4562}

\begin{longtable}[]{@{}
  >{\raggedright\arraybackslash}p{(\linewidth - 14\tabcolsep) * \real{0.1250}}
  >{\raggedright\arraybackslash}p{(\linewidth - 14\tabcolsep) * \real{0.1250}}
  >{\raggedright\arraybackslash}p{(\linewidth - 14\tabcolsep) * \real{0.1250}}
  >{\raggedright\arraybackslash}p{(\linewidth - 14\tabcolsep) * \real{0.1250}}
  >{\raggedright\arraybackslash}p{(\linewidth - 14\tabcolsep) * \real{0.1250}}
  >{\raggedright\arraybackslash}p{(\linewidth - 14\tabcolsep) * \real{0.1250}}
  >{\raggedright\arraybackslash}p{(\linewidth - 14\tabcolsep) * \real{0.1250}}
  >{\raggedright\arraybackslash}p{(\linewidth - 14\tabcolsep) * \real{0.1250}}@{}}
\toprule\noalign{}
\begin{minipage}[b]{\linewidth}\raggedright
Outcome
\end{minipage} & \begin{minipage}[b]{\linewidth}\raggedright
Mean
\end{minipage} & \begin{minipage}[b]{\linewidth}\raggedright
Median
\end{minipage} & \begin{minipage}[b]{\linewidth}\raggedright
25th pct
\end{minipage} & \begin{minipage}[b]{\linewidth}\raggedright
75th pct
\end{minipage} & \begin{minipage}[b]{\linewidth}\raggedright
95th pct
\end{minipage} & \begin{minipage}[b]{\linewidth}\raggedright
\% zeros
\end{minipage} & \begin{minipage}[b]{\linewidth}\raggedright
Planned model
\end{minipage} \\
\midrule\noalign{}
\endhead
\bottomrule\noalign{}
\endlastfoot
Views (\texttt{play}) & 570,285 & 7,716 & 239 & 245,933 & 3,140,807 &
1.3 & Negative binomial \\
Likes (\texttt{like}) & 19,265 & 109 & 5 & 4,694 & 105,566 & 5.1 &
Negative binomial \\
Favorites & 8,819 & 66 & 2 & 1,806 & 37,505 & 16.6 & Negative
binomial \\
Shares & 3,171 & 18 & 0 & 622 & 12,046 & 27.0 & \textbf{ZINB} \\
Comments (\texttt{review}) & 674 & 11 & 0 & 358 & 3,486 & 30.5 & ZINB \\
Coins & 4,682 & 7 & 0 & 282 & 14,663 & 36.5 & \textbf{ZINB} \\
Danmaku & 1,170 & 1 & 0 & 124 & 4,149 & 45.8 & \textbf{ZINB} \\
\end{longtable}

\textbf{Notes}: N = 4,562 unique videos. Distributions are heavily
right-skewed; we report median, 25th, 75th, and 95th percentiles in
addition to the mean to convey the dispersion. The ``\% zeros'' column
reports the percentage of videos with the corresponding outcome equal to
zero. ZINB = zero-inflated negative binomial. The high-cost endorsement
outcomes (coins, shares, danmaku) all exhibit zero-inflation rates above
25 percent, motivating the ZINB specification. The inflation equation in
the ZINB specification includes \texttt{log\_followers} and
\texttt{video\_age\_days}.

\begin{center}\rule{0.5\linewidth}{0.5pt}\end{center}

{[}Figure 1 about here{]}

\textbf{Figure 1 (planned)}: Density plot of the multimodal clickbait
index (factor score derived from emotion intensity + threat imagery +
text overlay intensity) on the 500-thumbnail human gold-standard
subsample, overlaid by a parallel distribution from the full-corpus VLM
coding, color-coded by Spearman ρ convergent-validity tier.

{[}Figure 2 about here{]}

\textbf{Figure 2 (planned)}: Schematic of the dual-track analytic
strategy. Left panel: the between-uploader cross-sectional specification
on the full 4,562-video sample. Right panel: the within-uploader
fixed-effects specification on the ≥ 2-video subsample (n ≈ 1,400
uploaders, ≈ 3,500 video-observations).

{[}Figure 3 about here{]}

\textbf{Figure 3 (planned)}: Hybrid annotation pipeline diagram. Tier 1:
500-thumbnail human gold standard (double-blind, two coders). Tier 2:
full-corpus computational coding (RapidOCR + RetinaFace+FER;
Chinese-CLIP for representativeness; Claude Opus 4.7 for VLM-scored
constructs). Tier 3: convergent-validity audit (Spearman ρ ≥ 0.60
threshold for full-corpus promotion).

\begin{center}\rule{0.5\linewidth}{0.5pt}\end{center}

\protect\phantomsection\label{refs}
\begin{CSLReferences}{1}{0}
\bibitem[\citeproctext]{ref-al-ali2024extra}
Al-Ali, Mohammed Nahar, and Safa M. Hamzeh. 2024. {``Extra Cues Extra
Views: A Multimodal Detection of {Arabic} Clickbait Thumbnail
Verbo-Visual Cues.''} \emph{Discourse \& Communication} 18 (1): 3--27.
\url{https://doi.org/10.1177/17504813231192203}.

\bibitem[\citeproctext]{ref-araujo2020automated}
Araujo, Theo, Irina Lock, and Bob van de Velde. 2020. {``Automated
Visual Content Analysis ({AVCA}) in Communication Research: A Protocol
for Large Scale Image Classification with Pre-Trained Computer Vision
Models.''} \emph{Communication Methods and Measures} 14 (4): 239--65.
\url{https://doi.org/10.1080/19312458.2020.1810648}.

\bibitem[\citeproctext]{ref-arel-bundock2024how}
Arel-Bundock, Vincent, Noah Greifer, and Andrew Heiss. 2024. {``How to
Interpret Statistical Models Using {marginaleffects} for {R} and
{Python}.''} \emph{Journal of Statistical Software} 111 (9): 1--32.
\url{https://doi.org/10.18637/jss.v111.i09}.

\bibitem[\citeproctext]{ref-bilibili2024q3}
Bilibili Inc. 2024. {``Bilibili Inc. Reports Third Quarter 2024
Financial Results.''} Form 6-K filed with U.S. Securities and Exchange
Commission.
\url{https://www.sec.gov/cgi-bin/browse-edgar?action=getcompany/&CIK=0001723690}.

\bibitem[\citeproctext]{ref-bird2005signaling}
Bird, Rebecca Bliege, and Eric Alden Smith. 2005. {``Signaling Theory,
Strategic Interaction, and Symbolic Capital.''} \emph{Current
Anthropology} 46 (2): 221--48. \url{https://doi.org/10.1086/427115}.

\bibitem[\citeproctext]{ref-blom2015click}
Blom, Jonas Nygaard, and Kenneth Reinecke Hansen. 2015. {``Click Bait:
Forward-Reference as Lure in Online News Headlines.''} \emph{Journal of
Pragmatics} 76: 87--100.
\url{https://doi.org/10.1016/j.pragma.2014.11.010}.

\bibitem[\citeproctext]{ref-bravo2025global}
Bravo, Isaac, Katharina Prasse, Stefanie Walter, Saffron O'Neill, and
Margret Keuper. 2025. {``Global Dynamics of Climate Change Imagery:
Emotional and Engagement Effects Across Visual Frames on
{Twitter}/{X}.''} \emph{Science Communication}.
\url{https://doi.org/10.1177/10755470251401900}.

\bibitem[\citeproctext]{ref-cao2025relevancy}
Cao, Junyi, Xiaolin Li, and Liang Zhang. 2025. {``Is Relevancy
Everything? {A} Deep-Learning Approach to Understand the Effect of
Image-Text Congruence.''} \emph{Management Science} 71 (12).
\url{https://doi.org/10.1287/mnsc.2022.01896}.

\bibitem[\citeproctext]{ref-chen2025danmaku}
Chen, Jiawei. 2025. {``Understanding Danmaku and Comment Interactions
Through Content Features and Video Popularity.''} \emph{Procedia
Computer Science}.

\bibitem[\citeproctext]{ref-donath2007signals}
Donath, Judith. 2007. {``Signals in Social Supernets.''} \emph{Journal
of Computer-Mediated Communication} 13 (1): 231--51.
\url{https://doi.org/10.1111/j.1083-6101.2007.00394.x}.

\bibitem[\citeproctext]{ref-dong2025impact}
Dong, Wei, Yu Liu, Wenxuan Wang, Liling Jiang, and Yi Yi. 2025. {``The
Impact of Danmaku Ritual Types on User Digital Engagement in Video-Based
Social Media: The Moderating Role of Influencer Types and Domains.''}
\emph{Psychology \& Marketing}. \url{https://doi.org/10.1002/mar.70003}.

\bibitem[\citeproctext]{ref-edelmann2020computational}
Edelmann, Achim, Tom Wolff, Danielle Montagne, and Christopher A. Bail.
2020. {``Computational Social Science and Sociology.''} \emph{Annual
Review of Sociology} 46: 61--81.
\url{https://doi.org/10.1146/annurev-soc-121919-054621}.

\bibitem[\citeproctext]{ref-geise2025effects}
Geise, Stephanie, and Yi Xu. 2025. {``Effects of Visual Framing in
Multimodal Media Environments: A Systematic Review of Studies Between
1979 and 2023.''} \emph{Journalism \& Mass Communication Quarterly} 102
(3): 796--823. \url{https://doi.org/10.1177/10776990241257586}.

\bibitem[\citeproctext]{ref-guo2026impact}
Guo, Lu, Yezhu Wang, Pei Li, Yutong Wang, and Yazheng Li. 2026. {``The
Impact of Headline Characteristics on Clicks: A Case Study of a
{Chinese} Local Medium.''} \emph{Journalism Practice} 20 (4): 1427--55.
\url{https://doi.org/10.1080/17512786.2024.2411624}.

\bibitem[\citeproctext]{ref-heley2022missing}
Heley, Kathryn, Anna Gaysynsky, and Andy J. King. 2022. {``Missing the
Bigger Picture: The Need for More Research on Visual Health
Misinformation.''} \emph{Science Communication} 44 (4): 514--27.
\url{https://doi.org/10.1177/10755470221113833}.

\bibitem[\citeproctext]{ref-isaac2025thirty}
Isaac, Mathew S. 2025. {``Thirty Years of Persuasion Knowledge Research:
From Demonstrating Effects to Building Theory to Increasing
Applicability.''} \emph{Consumer Psychology Review}.
\url{https://doi.org/10.1002/arcp.1107}.

\bibitem[\citeproctext]{ref-jones2015health}
Jones, Christina L., Jakob D. Jensen, Courtney L. Scherr, Natasha R.
Brown, Katheryn Christy, and Jeremy Weaver. 2015. {``The {Health}
{Belief} {Model} as an Explanatory Framework in Communication Research:
Exploring Parallel, Serial, and Moderated Mediation.''} \emph{Health
Communication} 30 (6): 566--76.
\url{https://doi.org/10.1080/10410236.2013.873363}.

\bibitem[\citeproctext]{ref-keib2018picture}
Keib, Kate, Camila Espina, Yen-I Lee, Bartosz W. Wojdynski, Dongwon
Choi, and Hyejin Bang. 2018. {``Picture This: The Influence of
Emotionally Valenced Images on Attention, Selection, and Sharing of
Social Media News.''} \emph{Media Psychology} 21 (2): 202--21.
\url{https://doi.org/10.1080/15213269.2017.1378108}.

\bibitem[\citeproctext]{ref-khawar2025analyzing}
Khawar, Salman, and Mark Boukes. 2025. {``Analyzing Sensationalism in
News on {Twitter} ({X}): Clickbait Journalism by Legacy Vs.
Online-Native Outlets and the Consequences for User Engagement.''}
\emph{Digital Journalism} 13 (8): 1482--502.
\url{https://doi.org/10.1080/21670811.2024.2394764}.

\bibitem[\citeproctext]{ref-king2020advancing}
King, Andy J., and Allison J. Lazard. 2020. {``Advancing Visual Health
Communication Research to Improve Infodemic Response.''} \emph{Health
Communication} 35 (14): 1723--28.
\url{https://doi.org/10.1080/10410236.2020.1838094}.

\bibitem[\citeproctext]{ref-klein2020when}
Klein, Elizabeth G., Kristin Roberts, Jennifer Manganello, Rebecca
McAdams, and Lara McKenzie. 2020. {``When Social Media Images and
Messages Don't Match: Attention to Text Versus Imagery to Effectively
Convey Safety Information on Social Media.''} \emph{Journal of Health
Communication} 25 (11): 879--84.
\url{https://doi.org/10.1080/10810730.2020.1853282}.

\bibitem[\citeproctext]{ref-kuiken2017effective}
Kuiken, Jeffrey, Anne Schuth, Martijn Spitters, and Maarten Marx. 2017.
{``Effective Headlines of Newspaper Articles in a Digital
Environment.''} \emph{Digital Journalism} 5 (10): 1300--1314.
\url{https://doi.org/10.1080/21670811.2017.1279978}.

\bibitem[\citeproctext]{ref-lang2000limited}
Lang, Annie. 2000. {``The Limited Capacity Model of Mediated Message
Processing.''} \emph{Journal of Communication} 50 (1): 46--70.
\url{https://doi.org/10.1111/j.1460-2466.2000.tb02833.x}.

\bibitem[\citeproctext]{ref-li2020picture}
Li, Yiyi, and Ying Xie. 2020. {``Is a Picture Worth a Thousand Words?
{A}n Empirical Study of Image Content and Social Media Engagement.''}
\emph{Journal of Marketing Research} 57 (1): 1--19.
\url{https://doi.org/10.1177/0022243719881113}.

\bibitem[\citeproctext]{ref-limpijankit2025detecting}
Limpijankit, Marvin, and John R. Kender. 2025. {``Detecting Cultural
Differences in News Video Thumbnails via Computational Aesthetics.''}
\url{https://arxiv.org/abs/2505.21912}.

\bibitem[\citeproctext]{ref-liu2025understanding}
Liu, Anita K. C., and Ozan Kuru. 2025. {``Understanding the Effects of
Visual Misinformation: A Systematic Review of 10 Years (2014--2024).''}
\emph{Mass Communication and Society}.
\url{https://doi.org/10.1080/15205436.2025.2549716}.

\bibitem[\citeproctext]{ref-liu2023factors}
Liu, Shaoqiang, Jinghong Xu, Zi'an Zhao, and Xiaojun Li. 2023.
{``Factors Affecting Trust in {Chinese} Digital Journalism: Approach
Based on Folk Theories.''} \emph{Media and Communication} 11 (4):
355--66. \url{https://doi.org/10.17645/mac.v11i4.7169}.

\bibitem[\citeproctext]{ref-lu2021capturing}
Lu, Yingdan, and Jennifer Pan. 2021. {``Capturing Clicks: How the
{Chinese} Government Uses Clickbait to Compete for Visibility.''}
\emph{Political Communication} 38 (1--2): 23--54.
\url{https://doi.org/10.1080/10584609.2020.1765914}.

\bibitem[\citeproctext]{ref-lu2022pervasive}
---------. 2022. {``The Pervasive Presence of {Chinese} Government
Content on {Douyin} Trending Videos.''} \emph{Computational
Communication Research} 4 (1): 68--97.
\url{https://doi.org/10.5117/CCR2022.2.002.LU}.

\bibitem[\citeproctext]{ref-lu2023unpacking}
Lu, Yingdan, and Cuihua (Cindy) Shen. 2023. {``Unpacking Multimodal
Fact-Checking: Features and Engagement of Fact-Checking Videos on
{Chinese} {TikTok} ({Douyin}).''} \emph{Social Media + Society} 9 (1).
\url{https://doi.org/10.1177/20563051231157591}.

\bibitem[\citeproctext]{ref-metzger2010social}
Metzger, Miriam J., Andrew J. Flanagin, and Ryan B. Medders. 2010.
{``Social and Heuristic Approaches to Credibility Evaluation Online.''}
\emph{Journal of Communication} 60 (3): 413--39.
\url{https://doi.org/10.1111/j.1460-2466.2010.01488.x}.

\bibitem[\citeproctext]{ref-molyneux2020aggregation}
Molyneux, Logan, and Mark Coddington. 2020. {``Aggregation, Clickbait
and Their Effect on Perceptions of Journalistic Credibility and
Quality.''} \emph{Journalism Practice} 14 (4): 429--46.
\url{https://doi.org/10.1080/17512786.2019.1628658}.

\bibitem[\citeproctext]{ref-munger2020news}
Munger, Kevin. 2020. {``All the News That's Fit to Click: The Economics
of Clickbait Media.''} \emph{Political Communication} 37 (3): 376--97.
\url{https://doi.org/10.1080/10584609.2019.1687626}.

\bibitem[\citeproctext]{ref-naveed2025thumbnailtruth}
Naveed, Wajiha, Zartash Afzal Uzmi, and Zafar Ayyub Qazi. 2025.
{``{ThumbnailTruth}: A Multi-Modal {LLM} Approach for Detecting
Misleading {YouTube} Thumbnails Across Diverse Cultural Settings.''}
\url{https://arxiv.org/abs/2509.04714}.

\bibitem[\citeproctext]{ref-oster2019unobservable}
Oster, Emily. 2019. {``Unobservable Selection and Coefficient Stability:
Theory and Evidence.''} \emph{Journal of Business \& Economic
Statistics} 37 (2): 187--204.
\url{https://doi.org/10.1080/07350015.2016.1227711}.

\bibitem[\citeproctext]{ref-peng2024automated}
Peng, Yilang, Irina Lock, and Albert Ali Salah. 2024. {``Automated
Visual Analysis for the Study of Social Media Effects: Opportunities,
Approaches, and Challenges.''} \emph{Communication Methods and Measures}
18 (2): 163--85. \url{https://doi.org/10.1080/19312458.2023.2277956}.

\bibitem[\citeproctext]{ref-powell2015clearer}
Powell, Thomas E., Hajo G. Boomgaarden, Knut De Swert, and Claes H. de
Vreese. 2015. {``A Clearer Picture: The Contribution of Visuals and Text
to Framing Effects.''} \emph{Journal of Communication} 65 (6):
997--1017. \url{https://doi.org/10.1111/jcom.12184}.

\bibitem[\citeproctext]{ref-scott2021you}
Scott, Kate. 2021. {``You Won't Believe What's in This Paper!
{C}lickbait, Relevance and the Curiosity Gap.''} \emph{Journal of
Pragmatics} 175: 53--66.
\url{https://doi.org/10.1016/j.pragma.2020.12.023}.

\bibitem[\citeproctext]{ref-shin2025emotion}
Shin, Jieun, Chris DeFelice, and Soojong Kim. 2025. {``Emotion Sells:
Rage Bait Vs.~Information Bait in Clickbait News Headlines on Social
Media.''} \emph{Digital Journalism} 13 (7): 1271--90.
\url{https://doi.org/10.1080/21670811.2025.2505566}.

\bibitem[\citeproctext]{ref-spence1973job}
Spence, Michael. 1973. {``Job Market Signaling.''} \emph{Quarterly
Journal of Economics} 87 (3): 355--74.
\url{https://doi.org/10.2307/1882010}.

\bibitem[\citeproctext]{ref-valkenburg2013differential}
Valkenburg, Patti M., and Jochen Peter. 2013. {``The Differential
Susceptibility to Media Effects Model.''} \emph{Journal of
Communication} 63 (2): 221--43.
\url{https://doi.org/10.1111/jcom.12024}.

\bibitem[\citeproctext]{ref-vanderweele2017sensitivity}
VanderWeele, Tyler J., and Peng Ding. 2017. {``Sensitivity Analysis in
Observational Research: Introducing the {E}-Value.''} \emph{Annals of
Internal Medicine} 167 (4): 268--74.
\url{https://doi.org/10.7326/M16-2607}.

\bibitem[\citeproctext]{ref-vultee2022clickbait}
Vultee, Fred, G. Scott Burgess, Darryl Frazier, and Kelsey Mesmer. 2022.
{``Here's What to Know about Clickbait: Effects of Image, Headline and
Editing on Audience Attitudes.''} \emph{Journalism Practice} 16 (1):
1--18. \url{https://doi.org/10.1080/17512786.2020.1793379}.

\bibitem[\citeproctext]{ref-wang2025decoding}
Wang, Yikai, Bin Hu, Chaolan Tang, and Xian Yang. 2025. {``Decoding
Clickbait: The Impact of Clickbait Types and Structures on Cognitive and
Emotional Responses in Online Interactions.''} \emph{Cyberpsychology,
Behavior \& Social Networking}.
\url{https://doi.org/10.1089/cyber.2024.0295}.

\bibitem[\citeproctext]{ref-xu2023clickbait}
Xu, Zhan, Mary Laffidy, and Lauren Ellis. 2023. {``Clickbait for Climate
Change: Comparing Emotions in Headlines and Full-Texts and Their
Engagement.''} \emph{Information, Communication \& Society} 26 (10):
1915--32. \url{https://doi.org/10.1080/1369118X.2022.2050416}.

\bibitem[\citeproctext]{ref-yang2022chinese}
Yang, An, Junshu Pan, Junyang Lin, Rui Men, Yichang Zhang, Jingren Zhou,
and Chang Zhou. 2022. {``{Chinese} {CLIP}: Contrastive Vision-Language
Pretraining in {Chinese}.''} \url{https://arxiv.org/abs/2211.01335}.

\bibitem[\citeproctext]{ref-yoon2024assessing}
Yoon, Yejun, Seunghyun Yoon, and Kunwoo Park. 2024. {``Assessing News
Thumbnail Representativeness: Counterfactual Text Can Enhance the
Cross-Modal Matching Ability.''} In \emph{Findings of the Association
for Computational Linguistics: ACL 2024}, 9009--24. Bangkok, Thailand.
\url{https://doi.org/10.18653/v1/2024.findings-acl.534}.

\bibitem[\citeproctext]{ref-zhang2019clicking}
Zhang, Xueying, and Shuhua Zhou. 2019. {``Clicking Health Risk Messages
on Social Media: Moderated Mediation Paths Through Perceived Threat,
Perceived Efficacy, and Fear Arousal.''} \emph{Health Communication} 34
(11): 1359--68. \url{https://doi.org/10.1080/10410236.2018.1489202}.

\end{CSLReferences}

\end{document}
